% ============================================================================
% APPENDIX RA: LIT-RABIN - Reference-Dependent Preferences & Behavioral Theory
% Evidence-Based Framework for Economic and Social Behavior
% ============================================================================
% Category: LIT- (Literature Integration)
% Language: English
% Code: RA
% Version: 1.0
% Last Updated: 2026-01-17
% Primary Dependencies: CORE-WHAT, CORE-AWARE, CORE-READY, LIT-KT
% EBF Integration: Formalizes reference-dependence, present bias, and belief formation
% ============================================================================

\chapter*{Appendix UNMAPPED_RA: Matthew Rabin Research Integration}
\addcontentsline{toc}{chapter}{Appendix UNMAPPED_RA: LIT-RABIN - Reference-Dependent Preferences \& Behavioral Theory}
\label{app:rabin}

% ============================================================================
% CROSS-REFERENCE MAP
% ============================================================================
\begin{tcolorbox}[colback=blue!5!white,colframe=blue!75!black,title=Cross-Reference Map]
\textbf{Builds On:}
\begin{itemize}
    \item Appendix U (LIT-KT): Prospect theory foundations
    \item Appendix UNMAPPED_C (CORE-WHAT): Utility dimensions
    \item Appendix UNMAPPED_AU (CORE-AWARE): Awareness mechanisms
\end{itemize}

\textbf{Supports:}
\begin{itemize}
    \item Appendix UNMAPPED_AV (CORE-READY): Willingness and action thresholds
    \item Appendix UNMAPPED_LO (LIT-LOEWENSTEIN): Visceral influences, projection bias
    \item Appendix R (LIT-THALER): Mental accounting, choice architecture
\end{itemize}

\textbf{Related Literature:}
\begin{itemize}
    \item Appendix K (LIT-FEHR): Social preferences formalization
    \item Appendix UNMAPPED_CA (LIT-CAMERER): Strategic behavior, game theory
    \item Appendix S (LIT-SUNSTEIN): Policy applications
\end{itemize}
\end{tcolorbox}

% ============================================================================
% CHAPTER LINKAGE
% ============================================================================
\begin{tcolorbox}[colback=green!5!white,colframe=green!75!black,title=Chapter Linkage]
\textbf{Primary:} Chapter 6 (Reference Structure) -- Reference-dependent utility formalization

\textbf{Secondary:}
\begin{itemize}
    \item Chapter 10 (Utility Architecture): FEPSDE dimensions with reference points
    \item Chapter 11 (Awareness): Self-awareness and naivete
    \item Chapter 12 (Readiness): Present bias and action thresholds
\end{itemize}
\end{tcolorbox}

% ============================================================================
% ABSTRACT
% ============================================================================
\begin{abstract}
This appendix integrates Matthew Rabin's foundational contributions to behavioral economics into the Evidence-Based Framework. Rabin's work provides rigorous theoretical foundations for reference-dependent preferences (with Kőszegi), present bias and self-control (with O'Donoghue), belief biases (law of small numbers), and social preferences (with Charness). We formalize six core axioms (UNMAPPED_RA-1 to UNMAPPED_RA-6) that connect his theoretical innovations to EBF. Key parameters include the projection bias coefficient ($\alpha_{proj} = 0.4$), present bias ($\beta = 0.7$), and loss aversion in reference-dependent utility ($\lambda = 2.0$). John Bates Clark Medal 2001; MacArthur Fellowship.
\end{abstract}

% ============================================================================
% QUICK REFERENCE
% ============================================================================
\begin{tcolorbox}[colback=gray!5!white,colframe=gray!75!black,title=Quick Reference]
\textbf{Appendix:} RA (LIT)

\textbf{Researcher:} Matthew Rabin (Harvard, Clark Medal 2001)

\textbf{Core Contributions:}
\begin{itemize}
    \item Reference-dependent preferences with expectations-based reference points
    \item Present bias: $\beta$-$\delta$ model with sophistication/naivete
    \item Belief biases: Law of small numbers, gamblers fallacy
    \item Social preferences: Fairness equilibria, reciprocity
\end{itemize}

\textbf{Key Parameters:}
\begin{itemize}
    \item Present bias: $\beta \approx 0.7$ (immediate discount)
    \item Projection bias: $\alpha_{proj} \approx 0.4$ (current state weight)
    \item Loss aversion: $\lambda \approx 2.0$ (in reference-dependent utility)
    \item Narrow bracketing: 50\% choose dominated when isolated
\end{itemize}

\textbf{EBF Parameters Updated:}
\begin{itemize}
    \item Reference point $r$ in $U(x|r)$ (expectations-based)
    \item $\beta$ (present bias parameter) in CORE-READY
    \item $\alpha_{proj}$ (projection bias) in CORE-AWARE
    \item Sophistication parameter $\hat{\beta}$ (beliefs about future self)
\end{itemize}
\end{tcolorbox}

% ============================================================================
% FUNDAMENTAL QUESTION
% ============================================================================
% -----------------------------------------------------------------------------
% APPENDIX SCOPE BOX
% -----------------------------------------------------------------------------

\begin{tcolorbox}[
    colback=orange!5!white,
    colframe=orange!75!black,
    title={\textbf{Appendix Scope: Ziel / In-Scope / Out-of-Scope / Constraints}},
    fonttitle=\bfseries
]

\textbf{Ziel (Objective):}
\begin{itemize}[nosep]
    \item How can we formally model the psychological regularities that cause systematic deviations from rational choice, and what are the welfare implications?
\end{itemize}

\textbf{In-Scope:}
\begin{itemize}[nosep]
    \item Fundamental Question
    \item Core Axioms: Rabin's Contributions
    \item Key Empirical Results
    \item Integration with Evidence-Based Framework
\end{itemize}

\textbf{Out-of-Scope (delegiert an andere Appendices):}
\begin{itemize}[nosep]
    \item LIT-KT $\rightarrow$ Appendix U
    \item CORE-WHAT $\rightarrow$ Appendix UNMAPPED_C
    \item CORE-AWARE $\rightarrow$ Appendix UNMAPPED_AU
    \item CORE-READY $\rightarrow$ Appendix UNMAPPED_AV
    \item LIT-LOEWENSTEIN $\rightarrow$ Appendix UNMAPPED_LO
\end{itemize}

\textbf{Constraints (Anwendungsgrenzen):}
\begin{itemize}[nosep]
    \item [TODO: Define when this appendix does NOT apply]
    \item [TODO: Define required prerequisites]
    \item [TODO: Define known limitations]
\end{itemize}

\textbf{Lieferobjekte (Deliverables):}
\begin{itemize}[nosep]
    \item 2 Table(s)
    \item 9 Equation(s)
    \item 12 Axiom(s)
    \item Worked Example(s)
\end{itemize}

\end{tcolorbox}


\section{Fundamental Question}

\begin{tcolorbox}[colback=red!5!white,colframe=red!75!black,title=Central Question]
\textbf{How can we formally model the psychological regularities that cause systematic deviations from rational choice, and what are the welfare implications?}
\end{tcolorbox}

Rabin's research program addresses the gap between psychology and economics:

\begin{enumerate}
    \item \textbf{Formalization}: Transform psychological insights into tractable economic models
    \item \textbf{Predictions}: Derive testable implications from behavioral assumptions
    \item \textbf{Welfare}: Develop frameworks for behavioral welfare economics
    \item \textbf{Policy}: Guide interventions that respect behavioral realities
\end{enumerate}

For EBF, Rabin provides the theoretical micro-foundations that underpin the framework's treatment of reference dependence, temporal preferences, and belief formation.

% ============================================================================
% THEORETICAL FOUNDATIONS
% ============================================================================
\section{Theoretical Foundations}

\subsection{Reference-Dependent Preferences (with Kőszegi)}

The Kőszegi-Rabin model formalizes Kahneman-Tversky's prospect theory with a crucial innovation: \textit{reference points are determined by rational expectations}.

\begin{equation}
U(c|r) = m(c) + n(c|r)
\end{equation}

where:
\begin{itemize}
    \item $m(c)$ = consumption utility (standard)
    \item $n(c|r)$ = gain-loss utility relative to reference point $r$
    \item $r$ = rational expectations about outcomes
\end{itemize}

\textbf{Key insight}: Reference point is not arbitrary---it's what you expected to get. This endogenizes the reference point and makes the model predictive.

\subsection{Present Bias and Self-Control (with O'Donoghue)}

The $\beta$-$\delta$ model captures hyperbolic discounting tractably:

\begin{equation}
U_t = u_t + \beta \sum_{\tau=1}^{T} \delta^\tau u_{t+\tau}
\end{equation}

where:
\begin{itemize}
    \item $\beta < 1$ = present bias (immediate gratification preference)
    \item $\delta$ = standard exponential discount factor
    \item $\beta = 1$ = time-consistent preferences
\end{itemize}

\textbf{Sophistication vs. Naivete}:
\begin{itemize}
    \item \textbf{Sophisticated}: Knows $\beta < 1$, may use commitment devices
    \item \textbf{Naive}: Believes future self has $\hat{\beta} = 1$, procrastinates
    \item \textbf{Partial naive}: $\beta < \hat{\beta} < 1$, underestimates bias
\end{itemize}

\subsection{Belief Biases: Law of Small Numbers}

Rabin's belief models explain seemingly contradictory biases through a unified framework:

\begin{equation}
P(\text{sequence is representative}) \propto \text{local representativeness}
\end{equation}

\textbf{Implications}:
\begin{itemize}
    \item \textbf{Gambler's fallacy}: After streak of heads, expect tails (small sample should look like population)
    \item \textbf{Hot hand fallacy}: After streak of successes, infer skill (small sample reveals true ability)
    \item Same underlying bias, different contexts
\end{itemize}

% ============================================================================
% CORE AXIOMS
% ============================================================================
\section{Core Axioms: Rabin's Contributions}

\subsection{Axiom UNMAPPED_RA-1: Expectations-Based Reference Points}

\begin{tcolorbox}[colback=blue!5!white,colframe=blue!75!black,title=Axiom UNMAPPED_RA-1: Expectations-Based Reference Points]
\textbf{Statement}: Reference points are determined by rational expectations about outcomes. People feel losses and gains relative to what they expected, not arbitrary anchors.

\textbf{Formal Definition}:
\begin{equation}
r = \mathbb{E}[c | I]
\end{equation}

where $I$ is the information set at the time of expectation formation.

\textbf{Implications}:
\begin{itemize}
    \item Attachment effect: Owning creates expectation of keeping $\rightarrow$ WTA > WTP
    \item News utility: Unexpected outcomes generate additional utility/disutility
    \item Endogenous risk attitudes: Risk aversion depends on expectations
\end{itemize}

\textbf{Evidence Tier}: 5 (Theoretical with experimental support)
\end{tcolorbox}

\subsection{Axiom UNMAPPED_RA-2: Present Bias with Varying Awareness}

\begin{tcolorbox}[colback=blue!5!white,colframe=blue!75!black,title=Axiom UNMAPPED_RA-2: Present Bias and Self-Awareness]
\textbf{Statement}: People exhibit present bias ($\beta < 1$), and their awareness of this bias ($\hat{\beta}$) affects outcomes. Partial naivete can be worse than full naivete or sophistication.

\textbf{Formal Definition}:
\begin{equation}
\text{Actual: } \beta \approx 0.7; \quad \text{Believed: } \hat{\beta} \in [\beta, 1]
\end{equation}

\textbf{Predictions}:
\begin{itemize}
    \item Sophisticated ($\hat{\beta} = \beta$): Uses commitment, may over-commit
    \item Naive ($\hat{\beta} = 1$): Procrastinates, surprised by own behavior
    \item Partial naive ($\beta < \hat{\beta} < 1$): Worst of both worlds
\end{itemize}

\textbf{Evidence Tier}: 2 (Experimental evidence on commitment demand)
\end{tcolorbox}

\subsection{Axiom UNMAPPED_RA-3: Projection Bias}

\begin{tcolorbox}[colback=blue!5!white,colframe=blue!75!black,title=Axiom UNMAPPED_RA-3: Projection Bias]
\textbf{Statement}: People project their current preferences onto their future selves, underestimating how much their tastes will change.

\textbf{Formal Definition}:
\begin{equation}
\tilde{u}(c, s', s) = (1-\alpha) \cdot u(c, s') + \alpha \cdot u(c, s)
\end{equation}

where $s$ = current state, $s'$ = future state, $\alpha \approx 0.4$ = projection bias.

\textbf{Implications}:
\begin{itemize}
    \item Shopping while hungry leads to over-purchasing
    \item Underestimate adaptation to chronic conditions
    \item Misprediction of future happiness
\end{itemize}

\textbf{Evidence Tier}: 2 (Field and lab experiments)
\end{tcolorbox}

\subsection{Axiom UNMAPPED_RA-4: Law of Small Numbers}

\begin{tcolorbox}[colback=blue!5!white,colframe=blue!75!black,title=Axiom UNMAPPED_RA-4: Law of Small Numbers]
\textbf{Statement}: People believe that small samples should be representative of the population, leading to systematic inference errors.

\textbf{Formal Definition}:
\begin{equation}
P(\text{HHHHH}|\text{fair coin}) < P(\text{HTHTH}|\text{fair coin}) \quad \text{(subjectively)}
\end{equation}

\textbf{Consequences}:
\begin{itemize}
    \item Gambler's fallacy: Expect reversals in random sequences
    \item Hot hand: Over-infer skill from short streaks
    \item Base rate neglect: Small samples trump prior probabilities
\end{itemize}

\textbf{Evidence Tier}: 1 (Extensive experimental evidence)
\end{tcolorbox}

\subsection{Axiom UNMAPPED_RA-5: Narrow Bracketing}

\begin{tcolorbox}[colback=red!5!white,colframe=red!75!black,title=Axiom UNMAPPED_RA-5: Narrow Bracketing]
\textbf{Statement}: People evaluate risky choices in isolation rather than as part of a portfolio, leading to rejection of attractive gamble combinations.

\textbf{Formal Definition}:
\begin{equation}
\text{Reject } g_1 \text{ AND Reject } g_2 \quad \text{even when} \quad \text{Accept } (g_1 + g_2)
\end{equation}

\textbf{Empirical Finding}:
\begin{itemize}
    \item 50\% of subjects choose dominated combinations when gambles presented separately
    \item Aggregation across decisions would often yield positive expected value
\end{itemize}

\textbf{Evidence Tier}: 1 (Experimental with Weizsäcker)
\end{tcolorbox}

\subsection{Axiom UNMAPPED_RA-6: Fairness in Games}

\begin{tcolorbox}[colback=green!5!white,colframe=green!75!black,title=Axiom UNMAPPED_RA-6: Reciprocal Fairness]
\textbf{Statement}: People care about others' intentions, not just outcomes. They reward kind actions and punish unkind ones, even at personal cost.

\textbf{Formal Definition}:
\begin{equation}
U_i = \pi_i + f_i \cdot f_j \cdot \pi_j
\end{equation}

where $f_i$ = $i$'s kindness to $j$, $f_j$ = $j$'s kindness to $i$.

\textbf{Implications}:
\begin{itemize}
    \item Fairness equilibria: Mutual kindness or mutual punishment
    \item Intention matters: Same outcome evaluated differently based on intent
    \item Strategic implications: Cooperation sustainable with reciprocators
\end{itemize}

\textbf{Evidence Tier}: 2 (Experimental games, complements Fehr-Schmidt)
\end{tcolorbox}

% ============================================================================
% KEY EMPIRICAL RESULTS
% ============================================================================
\section{Key Empirical Results}

\subsection{Parameter Estimates from Rabin's Research}

\begin{table}[h]
\centering
\caption{Key Parameters from Rabin's Research}
\label{tab:rabin-parameters}
\begin{tabular}{llcc}
\toprule
\textbf{Parameter} & \textbf{Context} & \textbf{Estimate} & \textbf{Source} \\
\midrule
Present bias $\beta$ & Various & $0.7$ & O'Donoghue-Rabin \\
Projection bias $\alpha$ & Prediction errors & $0.4$ & Loewenstein+ (2003) \\
Loss aversion $\lambda$ & Reference-dependent & $2.0$ & Kőszegi-Rabin \\
Narrow bracketing & Dominated choices & $50\%$ & Rabin-Weizsäcker \\
Naivete $\hat{\beta} - \beta$ & Self-prediction & $0.15$ & Various \\
\bottomrule
\end{tabular}
\end{table}

\subsection{The Calibration Theorem}

\begin{tcolorbox}[colback=yellow!5!white,colframe=yellow!75!black,title=Worked Example: Rabin's Calibration Theorem (2000)]
\textbf{Setup}: Expected utility maximizer who rejects 50-50 bet to win \$110 or lose \$100.

\textbf{Implication}: If this reflects diminishing marginal utility of wealth, then this person must:
\begin{itemize}
    \item Reject 50-50 bet to win \$1 billion or lose \$1,000
    \item At any wealth level!
\end{itemize}

\textbf{Conclusion}: Risk aversion over small stakes cannot be explained by diminishing marginal utility. Must be loss aversion or narrow bracketing.

\textbf{EBF Integration}: Supports reference-dependent utility over expected utility for small-stakes decisions. Risk attitudes are context-dependent ($\Psi$-modulated).
\end{tcolorbox}

% ============================================================================
% EBF INTEGRATION
% ============================================================================
\section{Integration with Evidence-Based Framework}

\subsection{Connection to 10C Framework}

\begin{table}[h]
\centering
\caption{Rabin's Contributions to 10C Framework}
\label{tab:rabin-9c}
\begin{tabular}{lll}
\toprule
\textbf{CORE} & \textbf{Rabin Contribution} & \textbf{Mechanism} \\
\midrule
CORE-WHAT (C) & Reference-dependent utility & $U(c|r) = m(c) + n(c|r)$ \\
CORE-WHEN (V) & Context shapes reference & $r = \mathbb{E}[c|I, \Psi]$ \\
CORE-AWARE (AU) & Projection bias, naivete & $\alpha_{proj}$, $\hat{\beta}$ \\
CORE-READY (AV) & Present bias & $\beta$-$\delta$ model \\
CORE-HOW (B) & Belief-preference interaction & $\gamma_{belief-choice} > 0$ \\
\bottomrule
\end{tabular}
\end{table}

\subsection{Theoretical Foundation for EBF}

Rabin's work provides the formal micro-foundations for several EBF components:

\begin{enumerate}
    \item \textbf{Reference Structure} (Chapter 6): Expectations-based reference points
    \item \textbf{Awareness Dimension} (Chapter 11): Sophistication vs. naivete
    \item \textbf{Readiness Dimension} (Chapter 12): Present bias creates gap between intention and action
    \item \textbf{Belief Formation}: Law of small numbers explains systematic prediction errors
\end{enumerate}

\subsection{Policy Implications}

\begin{tcolorbox}[colback=yellow!5!white,colframe=yellow!75!black,title=Policy Implications from Rabin]
\begin{enumerate}
    \item \textbf{Commitment Devices}: Help sophisticated present-biased agents
    \item \textbf{Sin Taxes}: Optimal when internalities exist (self-control failures)
    \item \textbf{Cooling-off Periods}: Address projection bias in purchases
    \item \textbf{Default Options}: Leverage narrow bracketing and inertia
    \item \textbf{Information Framing}: Account for expectations-based reference points
\end{enumerate}
\end{tcolorbox}

% ============================================================================
% CRITICAL FOUNDATIONS
% ============================================================================
\section{Critical Foundations}

\begin{tcolorbox}[colback=orange!5!white,colframe=orange!75!black,title=Critical Foundations]
\begin{enumerate}
    \item \textbf{Kahneman \& Tversky (1979)}: Prospect theory foundations
    \item \textbf{Laibson (1997)}: Hyperbolic discounting and golden eggs
    \item \textbf{Tversky \& Kahneman (1992)}: Cumulative prospect theory
    \item \textbf{Fehr \& Schmidt (1999)}: Inequity aversion (alternative social preference model)
    \item \textbf{Strotz (1956)}: Time inconsistency foundations
\end{enumerate}
\end{tcolorbox}

% ============================================================================
% OPEN ISSUES
% ============================================================================
\section{Open Issues and Research Frontier}

\begin{enumerate}
    \item \textbf{Reference Point Dynamics}: How do reference points adjust over time?
    \item \textbf{Heterogeneity}: Distribution of $\beta$, $\alpha$, $\hat{\beta}$ in population?
    \item \textbf{Welfare Measurement}: How to do welfare analysis with reference-dependent preferences?
    \item \textbf{Strategic Interaction}: Reference-dependence in games?
    \item \textbf{Neural Foundations}: Brain mechanisms underlying reference-dependence?
\end{enumerate}

% ============================================================================
% SUMMARY
% ============================================================================

% === AUTO-GENERATED ENTRIES (2026-02-22) ===

%%%%%%%%%%%%%%%%%%%%%%%%%%%%%%%%%%%%%%%%%%%%%%%%%%%%%%%%%%%%%%%%%%%%%%%%%%%%%%%
\subsubsection{2010: The Gambler''s and Hot-Hand Fallacies: Theory and Applicatio...}
%%%%%%%%%%%%%%%%%%%%%%%%%%%%%%%%%%%%%%%%%%%%%%%%%%%%%%%%%%%%%%%%%%%%%%%%%%%%%%%

\paragraph{Citation.}
Rabin, Matthew and Vayanos, Dimitri (2010). ``The Gambler''s and Hot-Hand Fallacies: Theory and Applications.'' \textit{Review of Economic Studies}, 77.

\paragraph{10C Integration.}
\begin{itemize}[nosep]
  \item \textbf{Domain}: [To be specified]
  \item \textbf{use\_for}: LIT-O, LIT-RABIN
\end{itemize}

% === END AUTO-GENERATED ENTRIES ===
\section{Summary}

\begin{tcolorbox}[colback=gray!5!white,colframe=gray!75!black,title=Chapter Summary]
Matthew Rabin's research provides EBF with rigorous theoretical foundations:

\textbf{Theoretical Contributions}:
\begin{itemize}
    \item Expectations-based reference points (with Kőszegi)
    \item $\beta$-$\delta$ model with sophistication/naivete (with O'Donoghue)
    \item Law of small numbers and belief biases
    \item Narrow bracketing and dominated choices
    \item Fairness equilibria and reciprocity
\end{itemize}

\textbf{Key Parameters}:
\begin{itemize}
    \item Present bias: $\beta \approx 0.7$
    \item Projection bias: $\alpha \approx 0.4$
    \item Loss aversion: $\lambda \approx 2.0$
    \item Narrow bracketing prevalence: $\sim 50\%$
\end{itemize}

\textbf{EBF Integration}:
\begin{itemize}
    \item Formalizes reference structure in CORE-WHAT
    \item Provides awareness parameters for CORE-AWARE
    \item Explains action gaps through present bias in CORE-READY
\end{itemize}

\textbf{Key Insight}: Psychological regularities can be formalized rigorously, generating testable predictions and policy guidance while maintaining welfare-relevant analysis.
\end{tcolorbox}

% ============================================================================
% GLOSSARY
% ============================================================================
\section{Glossary}

Key terms defined in this appendix (see also Appendix G for complete glossary):

\begin{description}
    \item[Present Bias ($\beta$)] Tendency to overweight immediate payoffs relative to future payoffs
    \item[Projection Bias] Tendency to project current preferences onto future selves
    \item[Reference Point] Benchmark against which outcomes are evaluated as gains or losses
    \item[Sophistication] Awareness of one's own present bias
    \item[Naivete] Underestimation of one's own present bias
    \item[Narrow Bracketing] Evaluating decisions in isolation rather than as portfolio
    \item[Law of Small Numbers] Belief that small samples should be representative
    \item[Fairness Equilibrium] Nash equilibrium where reciprocal fairness motivations are satisfied
\end{description}

% ============================================================================
% REFERENCES
% ============================================================================
\section{References}

\nocite{bcm_master}

\textbf{Primary Rabin Sources} (20 papers integrated):

\begin{itemize}
    \item Rabin, M. (1993). Incorporating Fairness into Game Theory. \textit{AER}.
    \item Rabin, M. (1998). Psychology and Economics. \textit{JEL}.
    \item O'Donoghue, T., \& Rabin, M. (1999). Doing It Now or Later. \textit{AER}.
    \item Rabin, M. (2000). Risk Aversion and Expected-Utility Theory: A Calibration Theorem. \textit{Econometrica}.
    \item O'Donoghue, T., \& Rabin, M. (2001). Choice and Procrastination. \textit{QJE}.
    \item Charness, G., \& Rabin, M. (2002). Understanding Social Preferences. \textit{QJE}.
    \item Rabin, M. (2002). Inference by Believers in the Law of Small Numbers. \textit{QJE}.
    \item Loewenstein, G., O'Donoghue, T., \& Rabin, M. (2003). Projection Bias. \textit{QJE}.
    \item O'Donoghue, T., \& Rabin, M. (2003). Self-Awareness and Self-Control.
    \item Eyster, E., \& Rabin, M. (2005). Cursed Equilibrium. \textit{Econometrica}.
    \item Kőszegi, B., \& Rabin, M. (2006). A Model of Reference-Dependent Preferences. \textit{QJE}.
    \item O'Donoghue, T., \& Rabin, M. (2006). Optimal Sin Taxes. \textit{JPubE}.
    \item Kőszegi, B., \& Rabin, M. (2007). Reference-Dependent Risk Attitudes. \textit{AER}.
    \item Kőszegi, B., \& Rabin, M. (2008). Choices, Situations, and Happiness. \textit{JPubE}.
    \item Kőszegi, B., \& Rabin, M. (2009). Reference-Dependent Consumption Plans. \textit{AER}.
    \item Rabin, M., \& Weizsäcker, G. (2009). Narrow Bracketing and Dominated Choices. \textit{AER}.
    \item Eyster, E., \& Rabin, M. (2010). Naive Herding. \textit{AEJ: Micro}.
    \item Rabin, M., \& Vayanos, D. (2010). Gambler's and Hot-Hand Fallacies. \textit{REStud}.
    \item Benjamin, D., Rabin, M., \& Raymond, C. (2016). Nonbelief in Law of Large Numbers. \textit{JEEA}.
\end{itemize}

\end{document}
